% In this file you should put the actual content of the blueprint.
% It will be used both by the web and the print version.
% It should *not* include the \begin{document}
%
% If you want to split the blueprint content into several files then
% the current file can be a simple sequence of \input. Otherwise It
% can start with a \section or \chapter for instance.

\section{The Artin--Schreier theorem}

\begin{lemma} \label{lem:p-power extensions}
\lean{X_pow_sub_C_irreducible_iff_of_prime_pow}
\mathlibok
Let $F$ be a field.
Let $p$ be an odd prime.
Then for $a \in F$, $a$ is not a $p$-th power
if and only if for every positive integer $n$, the polynomial $x^{p^n}-a$ is irreducible over $F$.
(Note that this fails for $p=2$, for example $x^4+4 = (x^2-2x+2)(x^2+2x+2)$.)
\end{lemma}
\begin{proof}
\mathlibok
In Mathlib.
\end{proof}

\begin{lemma} \label{lem:ac-quadratic-adjoins-i}
\lean{quadratic_algebraic_closure_no_i}
\leanok
Let $F$ be a field. Suppose that $F$ contains a square root $i$ of $-1$
and there exists a quadratic extension $K/F$ which is algebraically closed.
Then every element of $F$ is a square.
\end{lemma}
\begin{proof}
Let $a$ be an arbitrary element of $F$.
Since $K$ is algebraically closed, it contains some $b$ for which $b^4 = a$.
The minimal polynomial $f(x)$ of $b$ over $F$ has degree dividing $[K:F] = 2$.
If this degree is 1, then $f(x) = x-b$ and clearly $b^2 \in F$.
If the degree is 2, then $f(0)$ is the product of $b$ and another root of $f(x)$,
which must be one of $ib, -b, -bi$.
Hence $f(0) \in F$ equals $b^2$ times an element of $F$, so $b^2 \in F$.
In both cases, $b^2$ is a square root of $a$ in $F$.
\end{proof}

\begin{lemma} \label{lem:finite-ac-extension-separable}
\lean{finite_algebraic_closure_separable}
\leanok
Let $F$ be a field which admits a finite algebraically closed extension $K$.
Then $K/F$ is a separable extension.
\end{lemma}
\begin{proof}
\uses{lem:ac-quadratic-adjoins-i,lem:p-power extensions}
It suffices to check that $F$ is perfect; this is automatic if $F$ has characteristic $0$, so we
may assume that $F$ is of characteristic $p>0$, and then the perfect condition can be interpreted as the condition that every element of $F$ is a $p$-th power.

Suppose that $p=2$. Since $1 = -1$ in $F$, $F$ contains a square root of $-1$.
We can thus apply Lemma~\ref{lem:ac-quadratic-adjoins-i} to see that every element of $F$ is a perfect square, so $F$ is perfect.

Suppose that $p > 2$. If $F$ is not perfect, we can find $a \in F$ which is not a $p$-th power.
Choose a positive integer $n$ with $p^n > [K:F]$,
then choose $b \in K$ with $b^{p^n} = a$.
The polynomial $x^{p^n} - a$ is irreducible over $F$ by Lemma~\ref{lem:p-power extensions} and has $b$ as a root, so it is the minimal polynomial of $b$ over $F$.
However, the degree of the minimal polynomial is at most $[K:F] < p^n$, a contradiction.
\end{proof}


\begin{proposition} \label{prop:artin schreier extensions}
Let $F$ be a field of prime characteristic $p$.
Let $K/F$ be a Galois extension of degree $p$.
Then $K \cong F[x]/(x^p-x-a)$ for some $a \in F$.
\end{proposition}
\begin{proof}
To be filled in.
\end{proof}


\begin{lemma} \label{lem:cyclic-p-nonzero}
\lean{finite_algebraic_closure_cyclic_prime}
\leanok
Let $F$ be a field which admits an algebraically closed extension $K$ such that $K/F$ is Galois of prime degree $p$. Then $p \neq 0$ in $F$.
\end{lemma}
\begin{proof}
\uses{prop:artin schreier extensions}
If $p=0$ in $F$, then by Proposition~\ref{prop:artin schreier extensions},
$K = F[x]/(x^p-x-a)$ for some $a \in F$. Since $K$ is algebraically closed,
there exists $y \in K$ such that $y^p-y=ax^{p-1}$. Write $y = \sum_{i=0}^{p-1} y_i x^i$ with $y_i \in F$;
then $y_{p-1}^{p} - y_{p-1} = a$ and so $x-y_{p-1} \in \mathbb{F}_p$, a contradiction.
\end{proof}

\begin{lemma} \label{lem:ac-quadratic}
\lean{finite_algebraic_closure_cyclic_quadratic}
\leanok
Let $F$ be a field which admits an algebraically closed extension $K$ such that $K/F$ is Galois of prime degree $p$. Then $p = 2$.
\end{lemma}
\begin{proof}
\uses{lem:cyclic-p-nonzero,lem:p-power extensions}
Suppose by way of contradiction that $p \neq 2$.
By Lemma~\ref{lem:cyclic-p-nonzero}, $F$ cannot have characteristic $p$;
consequently, $K$ contains a primitive $p$-th root of unity $\zeta_p$.
Since $\Q(\zeta_p)/\Q$ is Galois of degree $p-1$, $F(\zeta_p)/F$ is Galois of degree dividing $p-1$;
but this degree must also divide $[K:F] = p$, so it is 1. That is, $\zeta_p \in F$.

By Kummer theory, $K  \cong F[x]/(x^p-a)$ for some $a \in F$. 
Lemma~\ref{lem:p-power extensions} implies that $x^{p^2}-a$ is irreducible over $F$,
so adjoining a root of this polynomial in $K$ gives a subfield of $K$ of degree $p^2$ over $F$.
However, the degree of any subfield must divide $[K:F] = p$, a contradiction.
\end{proof}


\begin{lemma} \label{lem:ac-finite-alg-closure-with-i}
\lean{finite_algebraic_closure_with_i}
\leanok
Let $F$ be a field containing a square root of $-1$ and admitting a finite algebraically closed extension. Then $F$ is algebraically closed.
\end{lemma}
\begin{proof}
\uses{lem:finite-ac-extension-separable,lem:ac-quadratic,lem:finite-ac-extension-separable,lem:ac-quadratic-adjoins-i}
The extension $K/F$ is normal, and by Lemma~\ref{lem:finite-ac-extension-separable} it is separable, so it is Galois; let $G$ be the Galois group. 
If $K \neq F$, then for any prime factor $p$ of $[L:K]$, by Cauchy's theorem 
$G$ contains a nontrivial cyclic subgroup of prime order $p$, whose fixed field is a subfield $E$ of $K$ containing $F$.
Applying Lemma~\ref{lem:ac-quadratic} to $K/E$ yields $p=2$.

By Lemma~\ref{lem:finite-ac-extension-separable}, $F$ is not of characteristic $2$.
We can thus write $K = F[x]/(x^2-a)$ for some $a \in F$ not a perfect square,
but by Lemma~\ref{lem:ac-quadratic-adjoins-i} no such $a$ exists.
\end{proof}

\begin{definition} \label{def:real closed field}
\lean{IsRealClosed}
\mathlibok
A field $F$ is \emph{real} if $-1$ is not a sum of squares in $F$.
A field $F$ is \emph{real closed} is all of the following conditions hold:
\begin{itemize}
\item $F$ is real;
\item for every $x \in F$, one of $x$ or $-x$ is a square in $F$;
\item every odd-degree polynomial over $F$ has a root in $F$.
\end{itemize}
\end{definition}

\begin{lemma} \label{lem:real closed condition}
\lean{RealClosed_from_quadratic}
\leanok
\uses{def:real closed field}
Let $F$ be a field. Suppose that $-1$ is not a square in $F$
and there exists a quadratic extension $K/F$ which is algebraically closed. Then $F$ is real closed.
\end{lemma}
\begin{proof}
Let $i \in K$ be a square root of $-1$.
The sum of any two squares in $F$ is a square: for $a,b \in F$, we can choose $u \in K$ with $u^2 = a+bi$, and then $\Norm_{K/F}(u)^2 = \Norm_{K/F}(a+bi) = a^2+b^2$.
Since $-1$ is not a square in $F$, it follows that $F$ is real.

If $x \in F$ is not a square, then $K=F(x^{1/2})$; by Kummer theory, $-1$ and $x$ belong to the same square class, and so $x^{1/2}/i \in F$ squares to $-x$.

Finally, since $F$ is real, it is of characteristic 0, so $K/F$ is a Galois extension.
For any polynomial $f(x)$ of odd degree over $F$, the Galois group of $K/F$ acts on the roots of $f$ with orbits of order 1 or 2, so there must be a fixed point; hence $f$ has a root in $F$.
\end{proof}

\begin{theorem} \label{thm:artin schreier thm}
\uses{def:real closed field}
\lean{artin_schreier_thm}
\leanok
Let $F$ be a field which admits a finite algebraically closed extension.
Then $F$ is algebraically closed or real closed.
\end{theorem}
\begin{proof}
\uses{lem:ac-finite-alg-closure-with-i,lem:real closed condition}
Let $K$ be an algebraic closure of $F$.
Let $i \in K$ be a square root of $-1$.
By Lemma~\ref{lem:ac-finite-alg-closure-with-i}, $F(i)$ is algebraically closed;
so either $i \in F$ and $F$ is algebraically closed, or $i \notin F$ and 
$F$ is real closed by Lemma~\ref{lem:real closed condition}.
\end{proof}


\section{Coherent sheaves and vector bundles}

Tag numbers refer to the Stacks Project.

\begin{definition}
\label{def:quasicoherent}
\mathlibok
(tag 01BE)
Let $X$ be a scheme. Let $F$ be a sheaf of $\cO_X$-modules.
We say that $F$ is \emph{quasicoherent} if for every $x \in X$,
there exist an open neighborhood $U$ of $x$ such that $\calF|_U$ is isomorphic to the kernel
of some morphism of $\cO_U$-modules of the form
\[
\bigoplus_{j \in J} \cO_U \to \bigoplus_{i \in I} \cO_U.
\]
\end{definition}

\begin{lemma}
Let $X$ be a scheme. Let $F$ be a sheaf of $\cO_X$-modules.
Then $F$ is quasicoherent if and only if 
\[
i^* F \to \tilde{M}, \qquad M := \Gamma(U, i^* F)
\]
of the tilde-Gamma adjunction is an isomorphism.
(Here $M$ is a module over the ring $R = \Gamma(X, \cO)$, and $X \cong \Spec R$.)
\end{lemma}
\begin{proof}
To be added to Mathlib.
\end{proof}

\begin{lemma}
Let $X$ be a scheme. Let $F$ be a sheaf of $\cO_X$-modules.
Then $F$ is quasicoherent if and only if for every ring $R$ and every morphism 
$i \colon U \to X$ with $U$ affine, the counit map
\[
i^* F \to \tilde{M}, \qquad M := \Gamma(U, i^* F)
\]
of the tilde-Gamma adjunction is an isomorphism.
\end{lemma}

\begin{definition}
Let $X$ be a scheme. Let $F$ be a sheaf of $\cO_X$-modules.
We say that $F$ is \emph{quasicoherent} if for every ring $R$ and every morphism 
$i \colon U \to X$ with $U$ affine, the counit map
\[
i^* F \to \tilde{M}, \qquad M := \Gamma(U, i^* F)
\]
of the tilde-Gamma adjunction is an isomorphism.
(Here $M$ is a module over the ring $\Gamma(U, \cO)$, and $U \cong \Spec R$.)
\end{definition}


\begin{definition}
Let $X$ be a scheme. Let $F$ be a sheaf of $\cO_X$-modules.
We say that $F$ is \emph{finitely generated} if $F$ is quasicoherent and,
for every ring $R$ and every morphism 
$i \colon U \to X$ with $U$ affine, the module $\Gamma(U, i^*F)$ over the ring $\Gamma(U, \cO)$ is finitely generated.
\end{definition}

\begin{lemma}
Let $X$ be a scheme. Let $F$ be a sheaf of $\cO_X$-modules.
Suppose that there exists an open covering of $F$ by affine subspaces $U_j$ such 
for each $j$, writing $i_j \colon U_j \to X$ for the inclusion, the module $\Gamma(U_j, i_j^*F)$ over the ring $\Gamma(U_j, \cO)$ is finitely generated. Then $F$ is finitely generated.
\end{lemma}
\begin{proof}
\end{proof}

\begin{definition} \label{def:vector bundles}
A \emph{vector bundle} on a scheme $X$ is a quasicoherent sheaf $F$ such that for every morphism 
$i \colon \Spec R \to X$, $\Gamma(X, i^* F)$ is a finite projective $R$-module. 
\end{definition}


\section{Abstract complete curves and vector bundles}

\begin{definition} \label{def:abstract curve}
By an \emph{abstract curve}, we mean a connected, separated, regular $1$-dimensional Noetherian scheme $X$. In other words, $X$ is a connected separated scheme which is covered by open subspaces which are the spectra of Dedekind domains.
\end{definition}

\begin{definition} \label{def:saturation}
\uses{def:abstract curve,def:vector bundles}
Let $F \subseteq G$ be an inclusion of vector bundles on an abstract curve $X$.
Then there is a unique intermediate bundle $F'$ such that $\rank(F) = \rank(F')$ and $G/F'$ is a vector bundle, called the \emph{saturation} of $F$ in $G$.
\end{definition}


\begin{definition} \label{def:abstract complete curve}
\uses{def:abstract curve}
By an \emph{abstract complete curve}, we mean an abstract complete curve
together with a \emph{degree} function from closed points of $X$ to positive integers with the property that the induced linear map $\deg\colon \Div(X) \to \Z$ is zero on all principal divisors.
\end{definition}


In the following lemmas, let $X$ denote an abstract complete curve.

\begin{definition} \label{def:degree map}
\uses{def:abstract complete curve,def:vector bundles}
Since the degree map $\Div(X) \to \Z$ vanishes on principal divisors, it induces a map $\Pic(X) \to \Z$
which we use to define a \emph{degree} homomorphism $\deg\colon K_0(X) \to \Z$.

If $F \subseteq F'$ is an inclusion of vector bundles of the same rank, then
$\deg(F) \leq \deg(F')$ with equality if and only if $F' = F$.
\end{definition}

\begin{definition} \label{def:complete flag}
\uses{def:abstract complete curve,def:vector bundles}
For $F$ a vector bundle on $X$, a \emph{complete flag} on $F$ is a filtration of the form
\[
0 = F_0 \subset \cdots \subset F_n = F
\]
in which $F_i/F_{i-1}$ is a vector bundle of rank $1$ for $i=1,\dots,n$.
Given a complete flag, define the \emph{slope sequence} as the tuple $(\deg(F_i/F_{i-1}))_{i=1}^n \in \Z^n$.
\end{definition}

\begin{lemma} \label{lem:filtration by line bundles}
\uses{def:complete flag}
Every vector bundle on $X$ admits a complete flag.
\end{lemma}
\begin{proof}
\uses{def:saturation}
We prove the theorem for a general vector bundle $F$ of rank $n$ by induction on $n$,
the case of rank 0 serving as a trivial base case.
For the induction step, pick a nonempty open affine subspace $U$ of $X$, let $Z$ be the reduced divisor $X \setminus U$, and choose a finite set of generators of $\Gamma(U, F)$ as a module over $\Gamma(U, \cO)$; these generators can be viewed as elements of $F(mZ)$ for some nonnegative integer $m$. We thus have a surjection $\cO_X^{\oplus d} \to F(mZ)$ for some positive integer $d$. In particular, there exists a nonzero (hence injective) map $\cO_X(mZ) \to F$;
choose one such map and let $F_1$ be the saturation of the image of this map. 
We then apply the induction hypothesis to $F/F_1$ to conclude.
\end{proof}

\begin{lemma} \label{lem:subbundle slope bound}
Let $F$ be a vector bundle on $X$ admitting a complete flag with degree sequence $m_1,\dots,m_n$.
Then for $i=0,\dots,n$, the degree of any rank-$i$ subbundle of $F$ is at most 
\[
\max\{m_{j_1} + \cdots + m_{j_i}\colon 1 \leq j_1,\dots,j_i \leq n\}.
\]
\end{lemma}
\begin{proof}
We prove the claim by strong induction on $n+i$.
By hypothesis, there exists a subbundle $F_1 \subseteq F$ of degree $m_1$
such that $F_1/F$ admits a complete flag with degree sequence $m_2,\dots,m_n$.
Let $G$ be an arbitrary subbundle of $F$ of rank $i$; we then have an exact sequence
\[
0 \to G \cap F_1 \to G \to G/(G \cap F_1) \to 0.
\]
If $G \cap F_1$ is zero,
then $G \cong/(G \cap F_1)$ is isomorphic to a subbundle of $F/F_1$ of rank $i$
and hence has degree at most
$\max\{m_{j_1} + \cdots + m_{j_i}: 2 \leq j_1,\dots,j_i \leq n\}$.
nonzero, then it is a line bundle of degree at most $m_1$,
while $G/(G \cap F_1)$ is a subbundle of $F/F_1$ of rank $i-1$ and hence has degree at most
$\max\{m_{j_1} + \cdots + m_{j_{i-1}}: 2 \leq j_1,\dots,j_{i-1} \leq n\}$.
In both cases,
\[
\deg(G) = \deg(G \cap F_1) + \deg(G/(G \cap F_1))
\]
satisfies the indicated bound.
\end{proof}

\begin{definition} \label{def:slope of vector bundle}
\uses{def:abstract complete curve,def:vector bundles}
Let $F$ be a nonzero vector bundle on $X$.
We define the \emph{slope} of $F$ as 
\[
\mu(F) := \frac{\deg(F)}{\rank(F)}.
\]
\end{definition}

\begin{definition} \label{def:stable bundle}
\uses{def:slope of vector bundle}
Let $F$ be a nonzero vector bundle on $X$.
The bundle $F$ is \emph{stable} if there is no nonzero proper subbundle $G$ of $F$ with $\mu(G) \geq \mu(F)$.
The bundle $F$ is \emph{semistable} if there is no nonzero subbundle $G$ of $F$ with $\mu(G) > \mu(F)$.
The bundle $F$ is \emph{polystable} if it is a direct sum of stable vector bundles, all of the same slope.
\end{definition}

\begin{lemma} \label{lem:map of stable}
\uses{def:stable bundle}
Let $F,G$ be two stable vector bundles on $X$ of the same slope $\mu$.
Then any nonzero morphism $f\colon F \to G$ is an isomorphism.
\end{lemma}
\begin{proof}
\uses{def:saturation}
Let $H \neq 0$ be the image of $f$.
Let $H'$ be the saturation of $H$ in $G$. 
Then
\[
\mu = \mu(F) \leq \mu(H) \leq \mu(H') \leq \mu(G) = \mu
\]
so all of these inequalities must be equalities. Since $G$ is stable this means that $H = H' = G$, and since $F$ is stable the kernel of $f$ must be zero.
\end{proof}

\begin{corollary} \label{cor:function field of curve}
The ring $H^0(X, \cO_X)$ is a field.
\end{corollary}
\begin{proof}
By Lemma~\ref{lem:map of stable}, every nonzero endomorphism of $\cO_X$ is invertible.
\end{proof}

\section{Generalized projective lines}

\begin{definition}
A \emph{generalized projective line} is an  abstract complete curve $X$ satisfying the following conditions.
\begin{enumerate}
\item[(a)] Every divisor of degree $0$ is principal. In other words the map $\Pic(X) \to \Z$
is an isomorphism.
\item[(b)] We have $H^1(X, \cO_X) = 0$. 
\item[(c)] There exists a closed point $x \in X$ of degree $1$.
\end{enumerate}
\end{definition}

\begin{lemma} \label{L:rank 2 extension}
Let $X$ be a generalized projective line.
Let $F$ be a vector bundle on $X$ admitting a complete flag with slope sequence
$-1,0,\dots,0,1$ (with $n$ zeroes for some nonnegative integer $n$).
Then $H^0(X, F) \neq 0$.
\end{lemma}
\begin{proof}
\uses{cor:function field of curve}
Let $E$ be the ring $H^0(X, \cO_X)$, which by Corollary~\ref{cor:function field of curve} is a curve.
By (c), we can choose a closed point $x \in X$ of degree $1$;
let $\kappa_x$ be the residue field of $x$.
By (a), the ideal sheaf of $x$ is isomorphic to $\cO_X(-1)$.
Define the bundle $H$ to fit into the commutative diagram
\[
\begin{tikzcd}
0 \arrow{r} & \cO_X(-1) \arrow{r} \arrow{d} & F \arrow{r}\arrow{d} & F/\cO_X(-1) \arrow{r} \arrow[d,equal] & 0 \\
0 \arrow{r} & \cO_X \arrow{r} & H \arrow{r} & F/\cO_X(-1) \arrow{r} & 0
\end{tikzcd}
\]
in which the left vertical arrow is injective with cokernel supported at $x$.
Since $\Ext^1_X(\cO_X, \cO_X) = H^1(X, \cO_X)$ vanishes by (b), $H$ admits a filtration with successive quotients $\cO_X^{\oplus (n+1)}, \cO_X(1)$.

We may assume that $H^0(X, H(-1)) = 0$, as otherwise $H$ contains the subbundle $\cO_X(1)$ whose intersection with $F$ has degree $\geq 0$ and hence
by (b) is either $\cO_X$ or $\cO_X(1)$.
Then $H^0(X, H)$ injects into $H^0(X, H/H(-1))$, which we reinterpret as the fiber $H_x$ of $H$ at $x$.
Using the exact sequence
\[
0 \to H^0(X, \cO_X^{\oplus n+1}) \to H^0(X, H) \to H^0(X, \cO_X(1)) \to H^1(X, \cO_X^{\oplus n}) = 0
\]
(where the last equality comes from (b)),
we express the image of $H^0(X, H) \to H_x$ 
as $E e_1 + \cdots + E e_{n+1} + \kappa_x e_{n+2}$
for some basis $e_1,\dots,e_{n+2}$ of $H_x$ over $\kappa_x$.

Meanwhile, the image of $H^0(X, F) \subset H^0(X, H)$ in $H_x$ is the intersection of the image of $H^0(X, H)$ with the image of $F_x \to H_x$; the latter is a codimension 1 subspace of $H_x$ and so can be written as the kernel of a nonzero $\kappa_x$-linear functional $\lambda\colon H_x \to \kappa_x$. 
Define $a \in E$, $b \in \kappa_x$ by
\[
(a,b) := \begin{cases} (0, 1) & \lambda(e_{n+2}) = 0 \\ (1, -\lambda(e_1)/\lambda(e_{n+2})) & \lambda(e_{n+2}) \neq 0; \end{cases}
\]
then $ae_1 + be_{n+2}$ is the image in $H_x$
of a nonzero element of $H^0(X, F)$.
\end{proof}

\begin{lemma} \label{lem:minimal filtration}
Let $X$ be a generalized projective line.
Let $F$ be a vector bundle of rank $n$ on $X$.
Then $F$ admits a complete flag whose slope sequence $m_1,\dots,m_n$ has the following properties.
\begin{enumerate}
\item[(a)]
For $i=1,\dots,n-1$, we have $m_{i+1} \leq m_{i}+1$.
\item[(b)]
There do not exist indices $1 \leq i < j \leq n$ such that for some $m \in \Z$,
\[
(m_i,\dots,m_j) = (m-1, m,\dots,m,m+1).
\]
\end{enumerate}
\end{lemma}
\begin{proof}
For each $i  \in \{1,\dots,n-1\}$, 
Lemma~\ref{lem:subbundle slope bound}
imposes a uniform upper bound on $m_1 + \cdots + m_i$,
and in all cases $m_1 +\cdots + m_n = \deg(F)$.
Consequently, we can choose a complete flag to minimize 
\[
\sum_{i=1}^n i m_i 
= (n+1) \deg(F) -\sum_{i=1}^{n-1} (m_1 + \cdots + m_i).
\]
We prove that this complete flag has the desired property.

Suppose by way of contradiction that (a) fails for some $i$. 
Then $F_{i+1}/F_{i-1}(m_i-1)$ admits a complete flag with degree sequence
$-1, m_{i+1}-m_i$; we then obtain a subbundle $(F_{i+1}/F_{i-1})(-m_i-1)$
admitting a complete flag with degree sequence $-1, 1$.
By Lemma~\ref{L:rank 2 extension}, $F_{i+1}/F_{i-1}(-m_i-1)$ admits $\cO_X$ as a (not necessarily saturated) subbundle; we thus obtain a new complete flag in which the terms
$m_i, m_{i+1}$ in the degree sequence are replaced by new terms $m_i', m_{i+1}'$ with $m_i' > m_i$ and $m_i' + m_{i+1}' = m_i + m_{i+1}$, contradicting minimality.

Suppose by way of contradiction that (b) fails for some $i,j$.
Then $(F_j/F_i)(-m)$ admits a complete flag with degree sequence $-1,0,\dots,0,1$.
By Lemma~\ref{L:rank 2 extension} again, 
$F_j/F_i(-m)$ admits $\cO_X$ as a (not necessarily saturated) subbundle;
we thus obtain a new complete flag in which the terms of $m_i, \dots, m_j$ in the degree sequence
are replaced by new terms $m_i', \dots, m_j'$ with $m_i' > m_i$
and $m_i' + \cdots + m_k' \geq m_i + \cdots + m_k$ for $k=i+1,\dots,j$,
again contradicting minimality.
\end{proof}

\begin{lemma} \label{lem:sequence argument1}
\leanok
Let $n$ be a positive integer.
Let $m_1,\dots,m_n$ be a sequence of integers satisfying 
the following conditions.
\begin{enumerate}
\item[(a)]
For $i=1,\dots,n-1$, we have $m_{i+1} \leq m_{i}+1$.
\item[(b)]
There do not exist indices $1 \leq i < j \leq n$ such that for some $m \in \Z$,
\[
(m_i,\dots,m_j) = (m-1, m,\dots,m,m+1).
\]
\end{enumerate}
Then for $1 \leq i < j \leq n$,
\[
m_j \leq m_i+1 \qquad (1 \leq i < j \leq n).
\]
\end{lemma}
\begin{proof}
\leanok
We prove the claim by induction on $j-i$.
The base case $j-i=1$ is covered by (a).
For the induction step, suppose by way of contradiction that $m_j> m_i+1$ for some $i,j$.
Set $m := m_i + 1$.
For $k = i+1, \dots, j-1$ we have $m_{k} \leq m_i+1$ and $m_j \leq m_{k}+1$ by the induction hypothesis;
combining with $m_j \geq m_i + 2$ shows that both of these must be equalities
and so $m_k = m, m_j = m+1$.
Now $(m_i, \dots, m_j) = (m-1, m, \dots, m, m+1)$, which contradicts (b).
\end{proof}


\begin{lemma} \label{lem:sequence argument2}
\leanok
Let $n$ be a positive integer.
Let $m_1,\dots,m_n$ be a sequence of integers satisfying 
the following conditions:
\begin{align}
m_1 +\cdots+m_n = 0 & \label{eq:bound2} \\
m_1 + \cdots + m_i \leq 0 & \quad (i=1,\dots,n-1)  \label{eq:bound1}\\
m_j \leq m_i+1 & \quad (1 \leq i < j \leq n). \label{eq:bound3}
\end{align}
Then $m_1 = \cdots = m_n = 0$.
\end{lemma}
\begin{proof}
\leanok
We proceed by induction on $n$, with trivial base case $n=1$.
For $n > 1$, note that \eqref{eq:bound1} and \eqref{eq:bound2} together imply $m_n \geq 0$.
On the other hand, \eqref{eq:bound1} implies $\min_{i<n} \{m_i\} \leq 0$, and then \eqref{eq:bound3}
implies $m_n \leq 1$. Finally, if $m_n = 1$, then \eqref{eq:bound3} implies that $m_i \geq 0$ for $i=1,\dots,{n-1}$, but then $m_1 + \cdots + m_n \geq 1$ in violation of \eqref{eq:bound2}.
We deduce that $m_n = 0$; now the induction hypothesis applies to $m_1,\dots,m_{n-1}$ to finish.
\end{proof}


\begin{theorem}
Let $X$ be a generalized projective line.
Then every semistable bundle of degree $0$ on $X$ is trivial.
\end{theorem}
\begin{proof}
Let $F$ be a semistable vector bundle of degree $0$ and rank $n$.
Choose a complete flag as in Lemma~\ref{lem:minimal filtration}.
By Lemma~\ref{lem:sequence argument1} and 
Lemma~\ref{lem:sequence argument2}, we have $m_1 = \cdots = m_n = 0$.
Since $\Ext^1_X(\cO_X, \cO_X) = H^1(X, \cO_X)$ vanishes by Lemma~\ref{lem:H1 of line bundles}, it follows that $F$ is trivial.
\end{proof}

\section{Harder--Narasimhan filtrations}

\begin{proposition} \label{prop:construct HN filtration}
\uses{def:stable bundle}
Every vector bundle $F$ on $X$ admits a unique filtration
\[
0 = F_0 \subset F_1 \subset \cdots \subset F_l = F
\]
in which each successive quotient $F_i/F_{i-1}$ is a semistable vector bundle of some slope $\mu_i$ and $\mu_1 > \cdots > \mu_l$.
\end{proposition}
\begin{proof}
We proceed by induction on $n = \rank(F)$.
Let $S \subseteq \Q$ be the set of slopes of nonzero subbundles of $F$.
The set $S$ is nonempty (it contains $\mu(F)$), bounded above
(by Lemma~\ref{lem:filtration by line bundles} and Lemma~\ref{lem:subbundle slope bound}),
and consists only of rational numbers with denominator bounded by $n$.
It thus admits an upper bound $\mu_1$; let $F_1$ be a subbundle of $F$ of $\mu_1$ of maximal rank
among such subbundles.
We then continue by induction.
%See for example \cite[Th\'eor\`eme~5.5.2]{FaFoCfv}.
\end{proof}

\begin{definition} \label{def:HN filtration}
\uses{prop:construct HN filtration}
For $F$ a vector bundle on $X$, the filtration on $F$ described in Proposition~\ref{prop:construct HN filtration}
is called the  \emph{Harder--Narasimhan filtration}, or \emph{HN filtration}, of $F$.
We define the multiset of \emph{HN slopes} of $F$ to consist of $\mu_i = \deg(F_i/F_{i-1})/\rank(F_i/F_{i-1})$ with multiplicity $\rank(F_i/F_{i-1})$.
\end{definition}

\begin{lemma} \label{lem:semistable from product}
\uses{def:stable bundle}
Let $F,G$ be two nonzero vector bundles on $X$. If $F \otimes G$ is semistable, then so are $F$ and $G$.
\end{lemma}
\begin{proof}
\uses{def:HN filtration}
Let $F_1, G_1$ be the first steps in the respective HN filtrations of $F$ and $G$;
then $\mu(F_1) \geq \mu(F)$, $\mu(G_1) \geq \mu(G)$. 
Then $F_1 \otimes G_1$ is a nonzero subbundle of $F \otimes G$ of slope $\mu(F_1)+\mu(G_1)$; since
$F \otimes G$ is semistable, we must have 
\[
\mu(F_1)+\mu(G_1) \leq \mu(F \otimes G) = \mu(F) + \mu(G) \leq \mu(F_1) + \mu(G_1).
\]
We conclude that $\mu(F_1) = \mu(F), \mu(G_1) = \mu(G)$ and so $F$ and $G$ are semistable.
\end{proof}


\section{Classification of vector bundles}

\begin{definition}
Let $P := \bigoplus_{d \geq 0} P_d$ be a graded ring satisfying the following conditions.

\begin{itemize}
\item[(i)]
The subring $P_0$ of $P$ is a field and $P_1 \neq 0$.
\item[(ii)]
The multiplicative monoid 
\[
\left( \bigcup_{n=0}^\infty P_n \setminus \{0\} \right)/P_0^\times
\]
is generated by $(P_1 \setminus \{0\})/P_{0}^{\times}$.
\item[(iii)]
For each nonzero $s \in P_1$, there exists a graded ring isomorphism 
\begin{equation} \label{eq:graded ring isom mod s}
P/sP \cong P_0 \oplus T L_s[T]
\end{equation}
for some field $L_s$ containing $P_0$. 
\end{itemize}

\end{definition}

\begin{proposition} \label{prop:abstract curve2}
The degree map $\deg\colon \Div(X) \to \Z$ induces an isomorphism $\Pic(X) \cong \Z$.
\end{proposition}
\begin{proof}
The map $\Pic(X) \to \Z$ carries $\cO_X(n)$ to $n$, so we need only check injectivity.
This follows from the fact that for any nonzero $s \in P_1$,
$P[s^{-1}]_0$ is a principal ideal domain.
\end{proof}

\begin{lemma} \label{lem:H1 of line bundles}
We have $H^i(X, \cO_X(n)) = 0$ for all $n \geq 0$.
\end{lemma}
\begin{proof}
For any $s,t \in P_1$ which are linearly independent over $P_0$, by \eqref{eq:graded ring isom mod s} we have $P_2 = P_1 s + P_1 t$.
By induction on $m+n+k$ we deduce that $P_{m+n+k} = P_{m+k} s^n + P_{n+k} t^m$ for any $m,n >0$ and $k \geq 0$, and hence
\[
H^1(X, \cO_X(n)) = \frac{P[(st)^{-1}]_n}{P[s^{-1}]_n + P[t^{-1}]_n} = 
\varinjlim_m \frac{P_{n+2m}}{P_{n+m}s^m + P_{n+m} t^m} = 0. 
\qedhere
\]
\end{proof}


\begin{lemma} \label{lem:filtration by line bundles om}
Every vector bundle $F$ on $X$ admits a filtration of the form
\[
0 = F_0 \subset \cdots \subset F_n = F
\]
in which each successive quotient $F_i/F_{i-1}$ is isomorphic to $\cO_X(m_i)$
for some $m_i \in \Z$.
\end{lemma}
\begin{proof}
Consider the filtration in Lemma~\ref{lem:filtration by line bundles}. By
Proposition~\ref{prop:abstract curve2}, each successive quotient is isomorphic to $\cO_X(m_i)$
for some $m_i \in \Z$.
\end{proof}

\begin{lemma} \label{L:rank 2 extension}
Let $F$ be a vector bundle on $X$ admitting a filtration with successive quotients
$\cO_X(-1), \cO_X^{\oplus n}, \cO_X(1)$ for some $n \geq 0$.
Then $H^0(X, F) \neq 0$.
\end{lemma}
\begin{proof}
Choose $s \in P_1$ nonzero and let $x \in |X|$ be the unique closed point at which $s$ vanishes. Define the bundle $H$ to fit into the commutative diagram
\[
\begin{tikzcd}
0 \arrow{r} & \cO_X(-1) \arrow{r} \arrow{d} & F \arrow{r}\arrow{d} & F/\cO_X(-1) \arrow{r} \arrow[d,equal] & 0 \\
0 \arrow{r} & \cO_X \arrow{r} & H \arrow{r} & F/\cO_X(-1) \arrow{r} & 0
\end{tikzcd}
\]
in which the left vertical arrow is injective with cokernel supported at $x$.
Since $\Ext^1_X(\cO_X, \cO_X) = H^1(X, \cO_X)$ vanishes by Lemma~\ref{lem:H1 of line bundles}, $H$ admits a filtration with successive quotients $\cO_X^{\oplus (n+1)}, \cO_X(1)$.

We may assume that $H^0(X, H(-1)) = 0$, as otherwise $H$ contains the subbundle $\cO_X(1)$ whose intersection with $F$ has degree $\geq 0$ and hence
by Proposition~\ref{prop:abstract curve2} is either $\cO_X$ or $\cO_X(1)$.
Then $H^0(X, H)$ injects into $H^0(X, H/H(-1))$, which we reinterpret as the fiber $H_x$ of $H$ at $x$.
Using the exact sequence
\[
0 \to H^0(X, \cO_X^{\oplus n+1}) \to H^0(X, H) \to H^0(X, \cO_X(1)) \to H^1(X, \cO_X^{\oplus n}) = 0
\]
(where the last equality comes from Proposition~\ref{lem:H1 of line bundles}),
we express the image of $H^0(X, H) \to H_x$ 
as $P_0 e_1 + \cdots + P_0 e_{n+1} + L_s e_{n+2}$
for some basis $e_1,\dots,e_{n+2}$ of $H_x$ over $L_s$.

Meanwhile, the image of $H^0(X, F) \subset H^0(X, H)$ in $H_x$ is the intersection of the image of $H^0(X, H)$ with the image of $F_x \to H_x$; the latter is a codimension 1 subspace of $H_x$ and so can be written as the kernel of a nonzero $L_s$-linear functional $\lambda\colon H_x \to L_s$. 
Define $a \in P_0$, $b \in L_s$ by
\[
(a,b) := \begin{cases} (0, 1) & \lambda(e_{n+2}) = 0 \\ (1, -\lambda(e_1)/\lambda(e_{n+2})) & \lambda(e_{n+2}) \neq 0; \end{cases}
\]
then $ae_1 + be_{n+2}$ is the image in $H_x$
of a nonzero element of $H^0(X, F)$.
\end{proof}


\begin{lemma} \label{lem:filtration by line bundles2}
With notation as in Lemma~\ref{lem:filtration by line bundles}, we can further ensure that 
\[
m_j \leq m_i+1 \qquad (1 \leq i < j \leq n).
\]
\end{lemma}
\begin{proof}
For each $i  \in \{1,\dots,n-1\}$, 
Lemma~\ref{lem:subbundle-slope-bound}
gives a uniform upper bound on $m_1 + \cdots + m_i$
over all filtrations as in Lemma~\ref{lem:filtration by line bundles},
while $m_1 +\cdots + m_n = \deg(F)$.
Consequently, we can choose such a filtration so as to minimize 
\[
\sum_{i=1}^n i m_i = 
= (n+1) \left( \sum_{i=1}^n m_i \right) -\sum_{i=1}^{n-1} (m_1 + \cdots + m_n).
\]
We prove that this filtration has the desired property by induction on $j-i$.

We start with the base case $j=i+1$.
Suppose by way of contradiction that this fails for some $i$. We then have an exact sequence
\[
0 \to \cO_X(m_i) \to F_{i+1}/F_{i-1} \to \cO_X(m_{i+1}) \to 0.
\]
By pulling back along an embedding $\cO_X(m_i+2) \to \cO_X(m_{i+1})$, we obtain a subbundle $G$ of $F_{i+1}/F_{i-1}$ and an exact sequence
\[
0 \to \cO_X(m_i) \to G \to \cO_X(m_i+2) \to 0.
\]
By Lemma~\ref{L:rank 2 extension}, $G$ admits $\cO_X(m_{i}+1)$ as a (not necessarily saturated) subbundle. We can thus replace the terms $m_i, m_{i+1}$ in our original sequence by new terms $m_i', m_{i+1}'$ with $m_i' > m_i$ and $m_i' + m_{i+1}' = m_i + m_{i+1}$, contradicting minimality.

We next proceed to the induction step $j-i> 1$.
Suppose by way of contradiction that this fails for some $i$.
By the induction hypothesis, we have $m_{j-1} \leq m_i+1$ and $m_j \leq m_{j-1}+1$,
so both of these must be equalities; in particular, $m_j = m_i + 2$.
For $k=i+1,\dots,j-1$, the induction hypothesis yields $m_i + 1 \leq m_k \leq m_j-1$,
so the sequence $(m_i, \dots, m_j)$ looks like $(m-1, m, \dots, m, m+1)$ for some $m$.
That is, $F_j/F_i$ admits a filtration with successive quotients $\cO_X(m-1), \cO_X(m)^{\oplus (j-i-1)}, \cO_X(m+1)$.
By Lemma~\ref{L:rank 2 extension} again, $F_j/F_i$ admits $\cO_X(m)$ as a (not necessarily saturated) subbundle, but again this would contradict minimality.
\end{proof}

\begin{lemma} \label{lem:sequence by bounds}
Let $n$ be a positive integer.
Let $m_1,\dots,m_n$ be a sequence of integers satisfying 
the following conditions:
\begin{align}
m_1 + \cdots + m_i \leq 0 & \quad (i=1,\dots,n-1)  \label{eq:bound1}\\
m_1 +\cdots+m_i = 0 & \quad (i=n) \label{eq:bound2} \\
m_j \leq m_i+1 & \quad (1 \leq i < j \leq n). \label{eq:bound3}
\end{align}
Then $m_1 = \cdots = m_n = 0$.
\end{lemma}
\begin{proof}
We proceed by induction on $n$, with trivial base case $n=1$.
For $n > 1$, note that \eqref{eq:bound1} and \eqref{eq:bound2} together imply $m_n \geq 0$.
On the other hand, \eqref{eq:bound1} implies $\max_i\{m_i\} \leq 0$, and then \eqref{eq:bound3}
implies $m_n \leq 1$. Finally, if $m_n = 1$, then \eqref{eq:bound3} implies that $m_i \geq 0$ for $i=1,\dots,{n-1}$, but then $m_1 + \cdots + m_n \geq 1$ in violation of \eqref{eq:bound2}.
We deduce that $m_n = 0$; now the induction hypothesis applies to $m_1,\dots,m_{n-1}$ to finish.
\end{proof}

\begin{proposition} \label{prop:semistable slope 0}
Every vector bundle on $X$ which is semistable of slope $0$ is trivial.
\end{proposition}
\begin{proof}
Let $F$ be a semistable vector bundle of degree $0$ and rank $n$,
and choose a filtration as in Lemma~\ref{lem:filtration by line bundles2}.
By Lemma~\ref{lem:sequence by bounds}, we have $m_1 = \cdots = m_n = 0$.
Since $\Ext^1_X(\cO_X, \cO_X) = H^1(X, \cO_X)$ vanishes by Lemma~\ref{lem:H1 of line bundles}, it follows that $F$ is trivial.
\end{proof}



\begin{lemma} \label{lem:one-step-modification}
Let $F$ be a vector bundle on $X$ with all HN slopes in $[0,1)$. Then 
for any closed point $x \in |X|$,
for $j_x \colon x \to X$ the canonical inclusion,
there exist a coherent sheaf $G$ on $x$ and an exact sequence
\[
0 \to \cO_X^{\oplus \rank(F)} \to F \to j_x^* G \to 0
\]
of coherent sheaves on $X$.
\end{lemma}
\begin{proof}
We show that for $i=1,\dots,\deg(F)$, we can find an exact sequence
\begin{equation} \label{eq:modification}
0 \to H_i \to F \to j_x^* G_i \to 0
\end{equation}
in which $H_i$ is a vector bundle of the same rank $F$ with degree $\deg(F)-i$ and all HN slopes nonnegative (and hence in $[0,1)$). 
For $i = \deg(F)$,
$H_i$ will then be semistable of degree $0$, and hence trivial by Proposition~\ref{prop:semistable slope 0}.

In the case $i=1$, it suffices to check the claim after replacing $F$ with its maximum quotient having all HN slopes in $(0,1)$; at this point, \emph{every} exact sequence of the form \eqref{eq:modification},
in which $H_1$ is a vector bundle of the same rank $F$ with degree $\deg(F)-1$, has the property that the HN slopes of $H_1$ are nonnegative.
We achieve the general case by applying this argument repeatedly.
\end{proof}

\begin{corollary} \label{cor:slopes to sections}
Let $F$ be a nonzero vector bundle on $X$ whose HN slopes are all nonnegative.
Then $F$ is generated by global sections and $H^1(X, F) = 0$.
\end{corollary}
\begin{proof}
By replacing $F$ with the first step of its HN filtration, we may reduce to the case where $F$ is semistable. By twisting, we may further ensure that $\mu(F) \in [0,1)$. We then deduce both claims from Lemma~\ref{lem:H1 of line bundles} and  Lemma~\ref{lem:one-step-modification}.
\end{proof}

\begin{lemma} \label{lem:end of semistable}
For any semistable vector bundle $F$ on $X$, $F^\vee \otimes F$ is trivial.
\end{lemma}
\begin{proof}
We immediately reduce to the case where $F$ is stable with slope in $[0,1)$. In this case, suppose by way of contradiction that $F^\vee \otimes F$ is not semistable, and let $G$ be the first step of its HN filtration. Since $G \otimes G$ has slope $2\mu(G) > \mu(G)$, 
the composition map $(F^\vee \otimes F) \otimes (F^\vee \otimes F) \to F^\vee \otimes F$
restricts to zero on $G \otimes G$.
By Corollary~\ref{cor:slopes to sections}, $G$ is generated by global sections;
consequently, $G \otimes G$ is generated by nonzero global sections of the form $s_1 \otimes s_2$.
Any such section corresponds to a pair of nonzero endomorphisms $F \to F$ which compose to zero; this contradicts Lemma~\ref{lem:map of stable}.
We deduce that $F^\vee \otimes F$ is semistable of slope $0$,
and hence trivial by Proposition~\ref{prop:semistable slope 0}.
\end{proof}

\begin{proposition} \label{prop:tensor product of semistable}
The tensor product of any two semistable vector bundles on $X$ is semistable.
\end{proposition}
\begin{proof}
Let $F,G$ be two semistable vector bundles on $X$.
By Lemma~\ref{lem:end of semistable}, 
\[
(F^\vee \otimes G^\vee) \otimes (F \otimes G)
\cong (F^\vee \otimes F) \otimes (G^\vee \otimes G)
\]
is semistable. By Lemma~\ref{lem:semistable-from-product}, $F \otimes G$ must be semistable.
\end{proof}

\begin{lemma} \label{cor:split sequence in order}
Any exact sequence of vector bundles on $X$ of the form
\[
0 \to F_1 \to F \to F_2 \to 0,
\]
in which every HN slope of $F_1$ is greater than or equal to every HN slope of $F_2$, is split.
\end{lemma}
\begin{proof}
By Proposition~\ref{prop:tensor product of semistable}, $F_2^\vee \otimes F_1$ has nonnegative HN slopes.
Now $\Ext^1_X(F_2, F_1) = H^1(X, F_2^\vee \otimes F_1)$ which vanishes by Corollary~\ref{cor:slopes to sections}.
\end{proof}

\begin{proposition} \label{prop:decompose as stables}
Every vector bundle on $X$ splits (nonuniquely) as a direct sum of stable bundles. In particular, the HN filtration always splits and every semistable bundle is \emph{polystable} (i.e., a direct sum of stable bundles of the same slope).
\end{proposition}
\begin{proof}
Both assertions follow directly from Lemma~\ref{cor:split sequence in order}.
\end{proof}

\begin{lemma} \label{lem:degree 1 stable2}
For $d$ a positive integer with $\dim_{P_0} P_1 > d$
and $c \in \Z$, there exists a semistable vector bundle on $X$ of rank $d$ and degree $c$.
\end{lemma}
\begin{proof}
We may assume without loss of generality that $d>1$ and $c \in \{1,\dots,d-1\}$.
Choose $s \in P_1$ nonzero and let $x \in |X|$ be the unique closed point at which $s$ vanishes.
Set $F := \cO_X^{\oplus d}$ and let $e_1,\dots,e_d$ be the images in $F_x$
of the standard basis vectors of $H^0(X, F) \cong P_0^{d}$.
The condition $\dim_{P_0} P_1 > d$ is equivalent to $\dim_{P_0} L_s \geq d$;
consequently, we can choose elements $f_1,\dots,f_c \in L_s$
such that $1,f_1,\dots,f_c$ are linearly independent over $P_0$.
Let $V \subset F_x$ be the $L_s$-span of $e_1 - f_1 e_d, \dots, e_c - f_d e_d$
and construct the subbundle $G$ of $F$
so that we have an exact sequence
\[
0 \to G \to F \to j_x^* V \to 0;
\]
then $G^\vee$ is semistable of degree $c$.
\end{proof}

\begin{lemma} \label{cor:unique stable}
For any $\lambda \in \Q$, there is at most one isomorphism class of stable vector bundles on $X$ of slope $\lambda$.
\end{lemma}
\begin{proof}
Let $F,G$ be two such bundles. By Proposition~\ref{prop:tensor product of semistable}, $F^\vee \otimes G$ is trivial. We can thus find a nonzero morphism $f\colon 
F \to G$, which must be an isomorphism by Lemma~\ref{lem:map of stable}.
\end{proof}
 

\begin{theorem} \label{thm:abstract classification}
Set
\[
S := \begin{cases} \Z & \mbox{in the classical case,} \\
\tfrac{1}{2} \Z & \mbox{in the twistor case,} \\
\Q & \mbox{otherwise}.
\end{cases}
\]
Then the following statements hold.
\begin{enumerate}
\item[(a)]
There exists exactly one isomorphism class of stable vector bundles $\cO_X(\lambda)$ of slope $\lambda$
if $\lambda \in S$ and none if $\lambda \notin S$.
\item[(b)]
For $\lambda \in S$ with denominator $d$ in lower terms, $\rank(\cO_X(\lambda)) = d$.
\item[(c)]
The endomorphisms of $\cO_X(\lambda)$ form a division algebra of degree $d$ over $P_0$.
\item[(d)]
For $\lambda, \lambda' \in S$, $\cO_X(\lambda) \otimes \cO_X(\lambda') \cong \cO_X(\lambda+\lambda')^{\oplus m}$
for some positive integer $m$.
\item[(d)]
Every vector bundle on $X$ splits (nonuniquely) as a direct sum of summands each of the form $\cO_X(\lambda)$ for some $\lambda \in S$.
\end{enumerate}
\end{theorem}
\begin{proof}
If $\dim_{P_0} P_1 = \infty$, then 
we may deduce (a) and (b) from Lemma~\ref{lem:degree 1 stable2} and Lemma~\ref{cor:unique stable}. Otherwise, by Theorem~\ref{thm:artin-schreier} we may argue as follows.
\begin{itemize}
\item 
In the classical case, $P = P_0[x,y]$
and $H^1(X, \cO_X(-1)) = 0$.
Hence any exact sequence of the form
\[
0 \to \cO_X \to F \to \cO_X(1) \to 0 
\]
splits, as then does any filtration as in Lemma~\ref{lem:filtration by line bundles2}, proving (a).
From this, (b) is obvious.
\item
In the twistor case, by identifying $P \otimes_{P_0} P_0(\sqrt{-1})$ with $P_0(\sqrt{-1})[x^2,xy,y^2]$ and applying the classical case, we see that $S \subseteq \tfrac{1}{2}\Z$.
We deduce (a) and (b) by applying Lemma~\ref{lem:degree 1 stable2} to construct a stable bundle of rank $2$ and degree $1$, then applying Lemma~\ref{cor:unique stable} for uniqueness.
\end{itemize}

Given (a)--(b), (c) follows from Lemma~\ref{lem:map of stable} and Lemma~\ref{lem:end of semistable}; 
(d) is immediate from Proposition~\ref{prop:tensor product of semistable} and Proposition~\ref{prop:decompose as stables};
and (e) is immediate from Proposition~\ref{prop:decompose as stables}.
\end{proof}
